% Standalone theorem insert for W33 papers.
% File: paper/w33_levi_cycle_quotient_exact_sequence_insert.tex

\section{The Levi Cycle Quotient Exact Sequence}
\label{sec:levi-cycle-quotient-exact-sequence}

Let $\mathcal L(W(3,3))$ be the Levi graph of $W(3,3)$: its vertices are the
points and lines of $W(3,3)$, and its edges are point-line incidences, i.e. flags.
Since $W(3,3)$ has $40$ points, $40$ lines, and $160$ flags, the Levi graph has
\[
        |V(\mathcal L)|=80,
        \qquad
        |E(\mathcal L)|=160.
\]
It is connected, so its cycle rank is
\[
        b_1(\mathcal L)=|E|-|V|+1=160-80+1=81.
\]

\begin{theorem}[Levi Cycle Quotient Exact Sequence]
Let $A\in\{-1,0,+1\}^{160\times1620}$ be the signed flag-quadrangle phase matrix.
The $79$-dimensional left kernel of $A$ is the signed cut/local-relation space of
the Levi graph, and the $81$-dimensional image is naturally identified with the
cycle space of the Levi graph.  Equivalently, there is an exact sequence
\[
        0\longrightarrow \operatorname{Cut}(\mathcal L)_{79}
        \longrightarrow \mathbb R^{\mathrm{Flags}}_{160}
        \longrightarrow H_1(\mathcal L;\mathbb R)_{81}
        \longrightarrow 0.
\]
The signed phase matrix realizes this quotient inside the quadrangle tight frame.
\end{theorem}

\begin{proof}[Finite certificate]
The certificate reconstructs $W(3,3)$ and its Levi graph.  It verifies connectedness,
so the cycle rank is $160-80+1=81$.  It also constructs the signed local point and
line four-flag relations, checks that their span has dimension $79$, and verifies that
this relation span equals $\ker(A^T)$.  Therefore the quotient of the $160$-dimensional
flag module by local point/line relations has dimension $160-79=81$, matching both
$\operatorname{rank}(A)$ and $b_1(\mathcal L)$.
\end{proof}

\begin{corollary}[Exact meaning of the protected $81$]
The protected $81$-dimensional phase image is not only spectrally detected by
$AA^T/160$.  It is exactly the Levi-graph cycle space of $W(3,3)$:
\[
        81=b_1(\mathcal L(W(3,3))).
\]
Thus
\[
        160=79+81
\]
is the standard cut/cycle decomposition of the point-line incidence graph.
\end{corollary}
